% =============================================================================
% cap6 §6.2 Contribuições
% =============================================================================
\section{Contribuições}
\label{sec:conclusao-contribuicoes}

A primeira contribuição é \textbf{arquitetural}: a especificação ponta-a-ponta de um pipeline em cascata que combina três bibliotecas pré-treinadas heterogêneas (camera trapping, re-ID de pets, re-ID de fauna selvagem) sob uma única convenção de entrada/saída, com tratamento explícito do mascaramento de pessoas exigido pela LGPD antes de qualquer persistência identificável. A segunda é \textbf{metodológica}: um protocolo de avaliação restrito a datasets públicos, com baseline mínimo por fase e janela sintética como sucedâneo controlado da operação de campo, que permite demonstrar viabilidade técnica antes de comprometer recursos com coleta sistemática. A terceira é \textbf{operacional}: a definição de uma tipologia única de \emph{kit} de instrumentação por baterias e câmeras IP comerciais, com matriz de pontos prevista no Apêndice~\ref{ap:matriz-pontos} e estrutura de scripts reproduzíveis no Apêndice~\ref{ap:pseudocodigos}, que reduzem a barreira de entrada da implementação para a equipe \emph{Gatosdoc2} sem exigir competência técnica especializada em \emph{deep learning} para a operação cotidiana. Em conjunto, as três contribuições convertem o problema --- até então tratado em planilhas e visitas pontuais --- em um sistema escalável e auditável, e estabelecem a base sobre a qual o impacto \textbf{social} (transparência da gestão da colônia, comunicação com a comunidade do Campus) e \textbf{de saúde pública} (detecção precoce de fauna sinantrópica relevante para zoonoses) pode ser construído de forma incremental.
